This appendix collects the technical calculations that were deliberately compressed in Sections~\ref{sec:ecstructure}, \ref{ch:differential_calculus}, and \ref{ch:holomorphic}. Its role is parallel to that of Appendix~\ref{app:model-acs}: it does not enlarge the program of the paper, but records the explicit computations behind the contact form, the horizontal differential $\delta$, and the first arithmetic holomorphic fields.

\subsection{Darboux reduction, the Reeb field, and the elementary Lie algebra}
\label{appD:darboux-reeb}

Recall the arithmetic contact form
\[
\alpha=da-(\mu\,du+\lambda a\,dv)
\]
on $M=\mathbb R^3_{(u,v,a)}$, together with the horizontal fields
\[
D_u=\partial_u+\mu\,\partial_a,
\qquad
D_v=\partial_v+\lambda a\,\partial_a.
\]
The following three short computations are used implicitly throughout Section~\ref{sec:ecstructure}.

\begin{proposition}
\label{prop:appD-darboux-reeb}
Assume $\mu\lambda\neq 0$.
Then:
\begin{enumerate}
    \item After the change of variables
    \[
    \tilde u=\mu u,
    \qquad
    \tilde v=\lambda v,
    \qquad
    x:=\tilde v,
    \qquad
    y:=a,
    \qquad
    z:=a-\tilde u,
    \]
    the contact form becomes
    \[
    \alpha = dz-y\,dx.
    \]

    \item The Reeb field of $\alpha$ is
    \[
    R=-\frac1\mu\,\partial_u.
    \]

    \item The Lie algebra generated by $D_u,D_v,\partial_a$ splits as
    \[
    \mathfrak g:=\mathrm{span}\{D_u,D_v,\partial_a\}
    \cong \mathbb R \oplus \mathfrak{aff}(1).
    \]
\end{enumerate}
\end{proposition}

\begin{proof}
Since
\[
 d\tilde u=\mu\,du,
 \qquad
 d\tilde v=\lambda\,dv,
 \qquad
 dz=da-d\tilde u,
\]
one obtains
\[
\alpha=da-d\tilde u-a\,d\tilde v=dz-y\,dx.
\]
This is the standard Darboux form.

For the Reeb field, write
\[
R=r_u\partial_u+r_v\partial_v+r_a\partial_a.
\]
Because
\[
d\alpha=-\lambda\,da\wedge dv,
\]
the condition $i_Rd\alpha=0$ gives $r_v=r_a=0$.
The normalization $\alpha(R)=1$ then becomes
\[
-\mu r_u=1,
\]
so $R=-\mu^{-1}\partial_u$.

Finally,
\[
[D_u,D_v]=\mu\lambda\,\partial_a,
\qquad
[D_u,\partial_a]=0,
\qquad
[D_v,\partial_a]=-\lambda\,\partial_a.
\]
Set
\[
C:=\partial_u=D_u-\mu\partial_a,
\qquad
H:=\lambda^{-1}D_v,
\qquad
E:=-\partial_a.
\]
Then $[C,H]=[C,E]=0$, while
\[
[H,E]=E.
\]
Thus $\mathrm{span}\{H,E\}\cong\mathfrak{aff}(1)$ and
\[
\mathfrak g=\mathbb RC\oplus\mathrm{span}\{H,E\}\cong\mathbb R\oplus\mathfrak{aff}(1).
\]
\end{proof}

This proposition is mainly organizational. It says that the arithmetic contact manifold is locally the standard contact model, but written in coordinates adapted to the add/multiply directions.

\subsection{Exact local commutator squares and infinitesimal density}
\label{appD:commutator-squares}

The main text records the local defect carried by the commutator of the two horizontal fields. For later reference it is useful to isolate the exact finite formulas.

\begin{proposition}
\label{prop:appD-commutator-square}
Let $\Phi_u^h$ and $\Phi_v^k$ be the time-$h$ and time-$k$ flows of $D_u$ and $D_v$. Then
\[
\Phi_u^h(u,v,a)=(u+h,v,a+\mu h),
\qquad
\Phi_v^k(u,v,a)=(u,v+k,e^{\lambda k}a).
\]
Define the two-step endpoints
\[
p_{uv}:=\Phi_v^k\!\circ\Phi_u^h(u,v,a),
\qquad
p_{vu}:=\Phi_u^h\!\circ\Phi_v^k(u,v,a).
\]
Then
\[
a(p_{uv})-a(p_{vu})=\mu h(e^{\lambda k}-1).
\]
If one closes the base rectangle by the inverse legs, the lifted commutator loop satisfies
\[
\bigl(\Phi_v^{-k}\circ\Phi_u^{-h}\circ\Phi_v^k\circ\Phi_u^h\bigr)(u,v,a)
=\bigl(u,v,a+\mu h(1-e^{-\lambda k})\bigr).
\]
\end{proposition}

\begin{proof}
The flow formulas follow by solving
\[
\dot u=1,\quad \dot v=0,\quad \dot a=\mu
\]
for $D_u$, and
\[
\dot u=0,\quad \dot v=1,\quad \dot a=\lambda a
\]
for $D_v$.
Therefore
\[
p_{uv}
=\Phi_v^k(u+h,v,a+\mu h)
=\bigl(u+h,v+k,e^{\lambda k}(a+\mu h)\bigr),
\]
whereas
\[
p_{vu}
=\Phi_u^h(u,v+k,e^{\lambda k}a)
=\bigl(u+h,v+k,e^{\lambda k}a+\mu h\bigr).
\]
Subtracting the assignment coordinates gives the first formula.

For the closed commutator loop, apply the four maps in order:
\[
(u,v,a)
\xmapsto{\Phi_u^h}(u+h,v,a+\mu h)
\xmapsto{\Phi_v^k}(u+h,v+k,e^{\lambda k}(a+\mu h)),
\]
then
\[
\xmapsto{\Phi_u^{-h}}(u,v+k,e^{\lambda k}(a+\mu h)-\mu h)
\xmapsto{\Phi_v^{-k}}(u,v,a+\mu h-\mu he^{-\lambda k}).
\]
This is exactly the claimed vertical drift.
\end{proof}

The two formulas encode two slightly different finite quantities:
\begin{itemize}
    \item the \emph{open-path torsion} between two histories with the same base endpoint,
    \item the \emph{closed-loop holonomy} after returning to the original $(u,v)$.
\end{itemize}
They are not equal at finite scale, but their first-order densities coincide.

\begin{corollary}
\label{cor:appD-infinitesimal-density}
One has
\[
\frac{a(p_{uv})-a(p_{vu})}{hk}=\mu\lambda+O(k),
\qquad
\frac{\Delta_a(h,k)}{hk}=\mu\lambda+O(k),
\]
where
\[
\Delta_a(h,k):=\mu h(1-e^{-\lambda k}).
\]
In particular,
\[
\lim_{h,k\to 0}\frac{a(p_{uv})-a(p_{vu})}{hk}
=
\lim_{h,k\to 0}\frac{\Delta_a(h,k)}{hk}
=\mu\lambda.
\]
\end{corollary}

\begin{proof}
Use the expansions
\[
e^{\lambda k}-1=\lambda k+O(k^2),
\qquad
1-e^{-\lambda k}=\lambda k+O(k^2).
\]
\end{proof}

This is the local bridge between the exact exponential weighting of the ACS and the constant curvature density seen by the contact structure and the $\delta$-calculus.

\subsection{The \texorpdfstring{$B_n$}{Bn} basis and the finite upward sweep}
\label{appD:Bn-basis-upward-sweep}

Section~\ref{ch:differential_calculus} introduces the scaled monomials
\[
\tilde v:=\lambda v,
\qquad
q:=e^{\tilde v},
\qquad
B_n(a,v):=e^{-n\tilde v}a^n=q^{-n}a^n.
\]
The point of the basis is that it linearizes the multiplicative direction and makes the $D_u$-calculus algorithmic.

\begin{lemma}
\label{lem:appD-Bn-actions}
For every $n\ge 0$ one has
\[
D_vB_n=0,
\qquad
D_uB_n=\mu n e^{-\tilde v}B_{n-1}=\mu n q^{-1}B_{n-1}.
\]
\end{lemma}

\begin{proof}
Since $B_n=q^{-n}a^n$, it is independent of $u$, so
\[
D_uB_n=\mu\partial_a(q^{-n}a^n)=\mu n q^{-n}a^{n-1}=\mu n q^{-1}B_{n-1}.
\]
For $D_v$, use $\partial_v q=\lambda q$ and $\partial_a a^n=na^{n-1}$:
\[
\partial_vB_n=-n\lambda q^{-n}a^n,
\qquad
\lambda a\partial_aB_n=\lambda n q^{-n}a^n.
\]
These two terms cancel.
\end{proof}

\begin{lemma}
\label{lem:appD-one-step-antiderivative}
If $P=P(v,q)$ is independent of $u$, then
\[
D_u\!\left(\frac{qP}{\mu(n+1)}B_{n+1}\right)=PB_n.
\]
Hence
\[
D_u^{-1}(PB_n)=\frac{qP}{\mu(n+1)}B_{n+1}
\]
within the class of fields independent of $u$.
\end{lemma}

\begin{proof}
Since $P$ is independent of $u$ and $a$, one has
\[
D_u\!\left(\frac{qP}{\mu(n+1)}B_{n+1}\right)
=\frac{qP}{\mu(n+1)}D_uB_{n+1}
=\frac{qP}{\mu(n+1)}\mu(n+1)q^{-1}B_n
=PB_n.
\]
\end{proof}

The genuinely useful step is what happens when the coefficient depends polynomially on $u$.

\begin{proposition}
\label{prop:appD-upward-sweep}
Let $P=P(u,v,q)$ be polynomial in $u$, and consider a target term $PB_n$.
Define the first lift
\[
G_0:=\frac{qP}{\mu(n+1)}B_{n+1}.
\]
Then
\[
D_uG_0=PB_n+R_1,
\qquad
R_1:=\frac{q\,\partial_uP}{\mu(n+1)}B_{n+1}.
\]
In particular, the $u$-degree of the remainder is one lower than that of $P$.
Iterating this correction therefore terminates after finitely many steps.
\end{proposition}

\begin{proof}
Use the product rule:
\[
D_uG_0
=\frac{q\,\partial_uP}{\mu(n+1)}B_{n+1}
+\frac{qP}{\mu(n+1)}D_uB_{n+1}.
\]
By Lemma~\ref{lem:appD-Bn-actions},
\[
D_uB_{n+1}=\mu(n+1)q^{-1}B_n,
\]
so the second term is exactly $PB_n$.
The first term is $R_1$.
Since $\partial_uP$ lowers the $u$-degree by one, repeated correction ends after finitely many iterations.
\end{proof}

\paragraph{A sample two-step sweep.}
Take $P(u,v,q)=u$.
The first lift is
\[
G_0=\frac{qu}{\mu(n+1)}B_{n+1},
\]
and Proposition~\ref{prop:appD-upward-sweep} gives
\[
D_uG_0=uB_n+\frac{q}{\mu(n+1)}B_{n+1}.
\]
Applying Lemma~\ref{lem:appD-one-step-antiderivative} to the remainder yields the corrected primitive
\[
G=\frac{qu}{\mu(n+1)}B_{n+1}-\frac{q^2}{\mu^2(n+1)(n+2)}B_{n+2},
\]
for which
\[
D_uG=uB_n.
\]
This is the finite ``upward sweep'' mechanism used implicitly in the main text.

\subsection{Factorization of the arithmetic Cauchy--Riemann operators}
\label{appD:factorization}

Let
\[
\partial_{\mathrm{AEG}}:=\frac12(D_u-iD_v),
\qquad
\bar\partial_{\mathrm{AEG}}:=\frac12(D_u+iD_v),
\]
acting on complex-valued fields on $M$.
For a field $F=f+ig$, the arithmetic Cauchy--Riemann equations are exactly
\[
\bar\partial_{\mathrm{AEG}}F=0.
\]
We now record the two operator factorizations that stand behind the second-order consequence in Section~\ref{ch:holomorphic}.

\begin{proposition}
\label{prop:appD-factorization}
Define the horizontal Laplacian by
\[
\Delta_H:=D_u^2+D_v^2.
\]
Then
\[
4\partial_{\mathrm{AEG}}\bar\partial_{\mathrm{AEG}}
=\Delta_H+i\mu\lambda\,\partial_a,
\qquad
4\bar\partial_{\mathrm{AEG}}\partial_{\mathrm{AEG}}
=\Delta_H-i\mu\lambda\,\partial_a.
\]
Consequently, every arithmetic holomorphic field satisfies
\[
(\Delta_H+i\mu\lambda\,\partial_a)F=0.
\]
\end{proposition}

\begin{proof}
Expand the first composition:
\[
4\partial_{\mathrm{AEG}}\bar\partial_{\mathrm{AEG}}
=(D_u-iD_v)(D_u+iD_v)
=D_u^2+D_v^2+i(D_uD_v-D_vD_u).
\]
Since
\[
[D_u,D_v]=\mu\lambda\,\partial_a,
\]
this becomes
\[
4\partial_{\mathrm{AEG}}\bar\partial_{\mathrm{AEG}}
=\Delta_H+i\mu\lambda\,\partial_a.
\]
The second identity is analogous.
If $\bar\partial_{\mathrm{AEG}}F=0$, then applying $\partial_{\mathrm{AEG}}$ gives
\[
\partial_{\mathrm{AEG}}\bar\partial_{\mathrm{AEG}}F=0,
\]
hence
\[
(\Delta_H+i\mu\lambda\,\partial_a)F=0.
\]
\end{proof}

The point of this formula is conceptual rather than merely formal. The curvature term does not appear as an external perturbation: it is forced by the non-commutation of the two horizontal arithmetic directions.

\begin{corollary}
\label{cor:appD-horizontal-jacobian}
Let $F=f+ig$ be arithmetic holomorphic, and define
\[
|F'|_H^2:=(D_uf)^2+(D_ug)^2=(D_vf)^2+(D_vg)^2.
\]
Then
\[
\delta f\wedge\delta g=|F'|_H^2\,du\wedge dv.
\]
\end{corollary}

\begin{proof}
Using
\[
\delta f=(D_uf)\,du+(D_vf)\,dv,
\qquad
\delta g=(D_ug)\,du+(D_vg)\,dv,
\]
one gets
\[
\delta f\wedge\delta g
=(D_ufD_vg-D_vfD_ug)\,du\wedge dv.
\]
The arithmetic Cauchy--Riemann equations imply
\[
D_vg=D_uf,
\qquad
D_vf=-D_ug,
\]
so the coefficient becomes
\[
(D_uf)^2+(D_ug)^2=|F'|_H^2.
\]
\end{proof}

\subsection{A genuinely arithmetic holomorphic coordinate and a rigidity check}
\label{appD:explicit-zeta}

The first examples in Section~\ref{ch:holomorphic} should not all come from classical holomorphic functions on the $(u,v)$-plane. The next calculation shows that there are already local fields with genuine $a$-dependence.

\begin{proposition}
\label{prop:appD-zeta}
Fix a simply connected open set on which a branch of the logarithm is chosen for
\[
\mu+i\lambda a.
\]
Define
\[
\zeta(u,a):=u+\frac{i}{\lambda}\log(\mu+i\lambda a).
\]
Then $\zeta$ is arithmetic holomorphic:
\[
\bar\partial_{\mathrm{AEG}}\zeta=0.
\]
\end{proposition}

\begin{proof}
Since $\zeta$ is independent of $v$, only the $a$-dependence matters.
Differentiate the logarithm:
\[
\partial_a\zeta
=\frac{i}{\lambda}\cdot\frac{i\lambda}{\mu+i\lambda a}
=-\frac{1}{\mu+i\lambda a}.
\]
Therefore
\[
D_u\zeta
=1+\mu\partial_a\zeta
=1-\frac{\mu}{\mu+i\lambda a}
=\frac{i\lambda a}{\mu+i\lambda a},
\]
and
\[
D_v\zeta
=\lambda a\partial_a\zeta
=-\frac{\lambda a}{\mu+i\lambda a}.
\]
Hence
\[
D_u\zeta+iD_v\zeta
=\frac{i\lambda a}{\mu+i\lambda a}-i\frac{\lambda a}{\mu+i\lambda a}=0,
\]
which is exactly $2\bar\partial_{\mathrm{AEG}}\zeta=0$.
\end{proof}

This field is already enough to show that arithmetic holomorphicity is not merely the old planar theory rewritten in new notation. The coordinate $u$ and the assignment variable $a$ are intertwined through the logarithm of the affine quantity $\mu+i\lambda a$.

A complementary one-line calculation explains the rigidity statement from the main text.

\begin{corollary}
\label{cor:appD-rigidity-a-only}
If a complex field $F$ depends only on $a$ and satisfies
\[
\bar\partial_{\mathrm{AEG}}F=0,
\]
then $F$ is constant.
\end{corollary}

\begin{proof}
If $F=F(a)$, then
\[
2\bar\partial_{\mathrm{AEG}}F
=(D_u+iD_v)F
=(\mu+i\lambda a)F'(a).
\]
For real $a$, the factor $\mu+i\lambda a$ never vanishes when $\mu\neq 0$, so $F'(a)=0$.
\end{proof}
